\minitoc 


% ===================================================================================================================
% Modèles réguliers

\section{Modèles réguliers}

Dans cette section, nous nous intéressons aux propriétés des modèles dit \emph{réguliers}, c'est à dire 
vérifiant plusieurs hypothèses de Cramer-Rao. 

\subsection{Hypothèses de Cramer-Rao}

\begin{proposition}
    Rappelons donc les hypothèses de Cramer-Rao. 
    \begin{itemize}
        \item[$(H_1)$] : Le modèle est \emph{identifiable} et \emph{dominé}. 
        $$ i.e \quad \mathcal{P}_\Theta << \mu, \quad \text{et} \quad \forall \theta \in \Theta, f_\theta = \frac{ \partial \myP_\theta}{ \partial \mu}, 
        \quad \text{et} \quad \forall \theta_1 \not = \theta_2, \; \mu(\{f_{\theta_1} = f_{\theta_2}\}) = 0 
        $$
        \item[$(H_2)$] : Le modèle est \emph{identifiable} et \emph{homogène}. 
        $$ 
            i.e \quad \forall \theta \in \Theta, \; \nabla f_\theta \in L^2( \mu) \quad \text{et} \quad 
            \forall \theta_1 \not = \theta_2, \; \mu(\{f_{\theta_1} = f_{\theta_2}\}) = 0 
        $$ 
        \item[$(H_3)$] : Pour tout $A \in \mathcal{A}$, la fonction $ \theta \in \Theta \longmapsto \myP_\theta(A)$, admet des 
            dérivées partielles telles que : 
            $$ 
                \frac{ \partial}{ \partial \theta_k} \myP_\theta(A) = \frac{ \partial}{ \partial \theta_k} \int_A f_\theta \; d \mu
                = \int_A \left( \frac{ \partial}{ \partial \theta_k} f_\theta \right) \;  d \mu 
            $$ 
        \item[$(H_4)$] : Pour tout $A \in \mathcal{A}$, la fonction $ \theta \in \Theta \longmapsto \myP_\theta(A)$, admet des 
            dérivées partielles d'ordre 2 telles que : 
            $$ 
                \frac{ \partial^2}{ \partial \theta_j \partial \theta_k} \myP_\theta(A) = \frac{ \partial^2}{ \partial \theta_j \partial \theta_k} \int_A f_\theta \; d \mu
                = \int_A \left( \frac{ \partial^2}{ \partial \theta_j \partial \theta_k} f_\theta \right) \;  d \mu 
            $$ 
    \end{itemize}
\end{proposition}

Les modèles dits \emph{réguliers} seront des modèles vérifiant ces hypothèses. Les calculs seront donc plus 
simples.

% \begin{example}[Modèle Régulier Binomial]
%     Posons $ \myP_\theta = \mathcal{B}(n, \theta)$ de densité par rapport à la mesure de comptage : 
%         $$ 
%             \forall x \in \mathcal{X} = \llbracket 1, n \rrbracket, \quad 
%             f_\theta(x) = \binom{n}{x} \theta^x (1 - \theta)^{n-x} 
%         $$ 
%     Dans le cas d'une seule dimension, où l'on considère un échantillon d'une variable aléatoire, on : 
%     \begin{align*}
%         \forall \theta \in \Theta, \quad f_\theta'(x) 
%         &= \binom{n}{x} x \theta^{x-1} (1 - \theta)^{n-x} \times (n-x) \theta^{x-1} (1 - \theta)^{n-x-1} \mathbbm{1}_{x \leqslant n} 
%     \end{align*}

% \end{example}

\begin{example}[Modèle Irrégulier]
    Soit un modèle $( \mathcal{X}, \mathcal{A}, \myP_\theta)$ de densité : 
        $$ f_\theta(x) = k_\theta g_\theta(x) \mathbbm{1}_{[0, \theta]} $$
    Le support de $ \myP_\theta$ varie en fonction du paramètre, il n'est donc pas dominé 
    pour tout $ \theta \in \Theta$. Il ne vérifie donc pas $(H_1)$. 
    D'autre part, $(H_2)$ n'est pas vérifiée, car $ x \longmapsto \mathbbm{1}_{[0, \theta]}$ n'est pas dérivable 
    partout à cause du saut en $ \theta$. 
\end{example}

\subsection{Rappel d'Outils}

Effectuons quelques rappels sur des théorèmes utiles. 

\begin{rappel}[Théorème d'inversion globale]
    S'il existe $ g \in L^1( \mu)$ et $ \mathcal{N} \in \mathcal{A}$ une partie négligeable vérifiants : 
    \begin{enumerate}[label=\roman*)]
        \item $ \nabla_\theta f_\theta(x) $ existe pour tout $ \theta \in \Theta$ avec $x \not \in \mathcal{N}$ fixé, 
        \item $ \displaystyle \left| \frac{ \partial f_\theta}{ \partial \theta_k}\right| \leqslant g $
            pour tout $ k \in \llbracket 1, p \rrbracket$ et $ \theta \in \Theta$, 
    \end{enumerate}
    alors $ \theta \longmapsto \myP_\theta(A)$ est différentiable sur $ \Theta$ et : 
        $$ 
            \forall A \in \mathcal{A}, \forall \theta \in \Theta, \quad 
            \nabla_\theta \myP_\theta(A) = \int_A \nabla_\theta f_\theta \; d \mu 
        $$
\end{rappel}

\begin{rappel}[Théorème de Dérivation Locale]
    Soit $ \theta_0 \in \Theta$, $ \mathcal{N}_{\theta_0} \in \mathcal{A}$, et $V_0 \subset \Theta$ un voisinnage de $ \theta_0$ 
    tels que $ \mu( \mathcal{N}_{\theta_0}) = 0$. Si 
    \begin{enumerate}[label=\roman*)]
        \item $ \nabla_\theta f_\theta(x)$ existe en $ \theta_0$ pour tout $ x \not \in \mathcal{N}_{\theta_0}$, 
        \item $ | f_\theta(x) - f_{\theta_0}(x) | \leqslant g_{\theta_0}(x) \| \theta - \theta_0 \|_p $
            pour tout $ \theta \in V_0$ et $g_{\theta_0} \in L^1( \mu)$, 
    \end{enumerate}
    alors pour tout $A \in \mathcal{A}$, $ \theta \longmapsto \myP_\theta(A)$ est \emph{différentiable} en $ \theta = \theta_0$ 
    et : 
        $$ \nabla_\theta \int_A f_\theta \; d \mu = \int_A \nabla_\theta f_\theta \; d \mu $$
\end{rappel}

% ===================================================================================================================
% Vecteurs Scores

\section{Vecteurs Scores}

\subsection{Score}

Durant tout ce chapitre, nous serons amenés à calculer le gradient de log-vraissemblance. 
Pour simplifier les calculs et étudier leurs propriétés, nous allons les définir sous 
la forme de \emph{vecteurs scores}. 

\begin{definition}[Vecteur Score]
    Soit $(H_{1,2})$, le vecteur aléatoire score d'un échantillon $X$ est défini par le gradient de la 
    log-vraissemblance de l'échantillon. On a donc : 
    \begin{equation}
        \boxed{S_\theta(X) = \nabla_\theta \log f_\theta(X) = \frac{1}{f_\theta(X)} \nabla f_\theta(X). }
    \end{equation}
    C'est une \emph{variable aléatoire}. 
\end{definition}

\begin{prop}[Espérance]
    Sous $(H_{1, 2, 3})$, les vecteurs scores donc centrés. 
        $$ i.e \quad (H_{1, 2, 3}) \Longrightarrow \forall \theta \in \Theta, \; \mathbb{E}_\theta(S_\theta(X)) = 0_p $$
\end{prop}

\begin{proof}
    Supposons $(H_{1, 2, 3})$, on a alors que $ \int_\mathcal{X} f_\theta \; d \mu = 1$ car c'est une densité. 
    Donc par $(H_{1, 2, 3})$, on a : 
        \begin{align*}
            & \int_\mathcal{X} \frac{ \partial}{ \partial \theta_k} f_\theta \; d \mu = 0 \\ 
            \iff & \int_\mathcal{X} \frac{1}{f_\theta(x)} \frac{\partial}{\partial \theta_k} f_\theta(x) \times 
                \underbrace{ f_\theta(x) \; d \mu(x)}_{d \myP_\theta(x)} = 0 \\ 
            \iff & \int_\mathcal{X} \frac{ \partial }{\partial \theta_k} \log f_\theta \; d \myP_\theta = 0 \\ 
            \iff & \mathbb{E}_\theta \left( \frac{\partial}{\partial \theta_k} \log f_\theta \right) = 0 
        \end{align*}
\end{proof}

\subsection{Modèles exponentiels}

Voyons à présent un modèle régulier très utilisé : le \emph{modèle exponentiel}. 

\begin{definition}[Modèle Exponentiel]
    Un modèle $( \mathcal{X}, \mathcal{A}, \mathcal{P}_\Theta)$ sera dit \emph{exponentiel} si : 
        $$ \boxed{\forall \theta \in \Theta, \quad f_\theta(x) = k_\theta h(x) \exp \left( \langle T(x), \gamma(\theta) \rangle\right).} $$
    Avec $ \gamma( \Theta) \subset \R^q$, $T : \mathcal{X} \longrightarrow \R^q$ et 
    $h : \mathcal{X} \longrightarrow \R$ la fonction de changement de mesure. 

    Il sera qualifié \emph{de plein rang} si $p=q$. On dira aussi qu'il est \emph{canonique} sir $ p=q$ et 
    $ \gamma$ est l'identité. On appellera enfin $T(X)$ la \emph{statistique privilégiée du modèle}. 

    Par soucis de clarté, on notera souvent $ \kappa_\theta = \log k_\theta$. 
\end{definition}

\begin{theorem}[Régularité]
    Un modèle exponentiel canonique de plein rang vérifie $(H_{1, 2, 3, 4})$ et $ k_\theta \in \mathcal{C}^\infty$.
\end{theorem}

\begin{proposition}
    Les modèles canoniques sont très pratiques à manipuler. On a :  
        $$ f_\theta(x) = k_\theta h(x) \exp \left( \langle T(x), \theta \rangle\right) $$ 
    Après calcul, on obtient que $S_\theta(X) = \nabla_\theta \log k_\theta + T(X) \in \R^p$. 
    Or, on sait que $ \mathbb{E}_\theta(S_\theta(X)) = 0$, donc : 
        $$ \mathbb{E}_\theta(T(X)) = - \nabla_\theta \kappa_\theta. $$
\end{proposition}

\begin{prop}[Cas général]
    Dans un cadre plus général, un modèle exponentiel de plein rang admet des scores de la forme : 
        $$ S_\theta(X) = J_{ \gamma}^t \times T(X) + \nabla_\theta \kappa_\theta $$ 
    avec $ \displaystyle J_{\gamma} = \left( \left( \frac{\partial \gamma_i}{ \partial \theta_j} (\theta)\right)_{i,j}\right)$ 
    ma matrice Jacobienne $ p \times p$ de $ \gamma : \Theta \longrightarrow \R^p$. 
    Si $J_{\gamma}$ est définie positive, alors : 
        $$ \mathbb{E}_\theta(T(X)) = - \left( J_\gamma ^{-1}\right)^t \times \nabla_\theta \kappa_\theta $$
\end{prop}

% ===================================================================================================================
% Information de Fischer

\section{Information de Fischer}

Passons à la définition principale de ce chapitre : l'information de Fisher. Elle quantifie 
l'information contenue sur $ \theta$ dans l'acte d'observer $X$. 

\begin{definition}[Information de Fisher]
    Sous $(H_{1,2})$, la \emph{quantité d'information de Fisher} en $ \theta$ apportée par l'acte d'observer 
    une réalisation $X( \omega)$ sous $ \myP_\theta$ est : 
    \begin{equation}
        \mathcal{I}_\theta(X) = \mathbb{E}_\theta(S_\theta(X) \times S_\theta^t(X))
    \end{equation}
    Si $(H_4)$ est vérifiée, alors $ \mathcal{I}_\theta(X) = \V_\theta(S_\theta(X))$, la 
    matrice de covariance du vecteur score. 
\end{definition}

\begin{prop}[Information de Fisher]
    \begin{itemize}
        \item Sous $(H_{1,2})$, $ \mathcal{I}_\theta(X)$ est intrinsèque au modèle (i.e elle ne dépend plus de 
        la mesure dominante $ \mu$). 
        \item Sous $(H_{1,2,3,4})$, on a : 
        $$ 
            \mathcal{I}_\theta(X) = - \mathbb{E}_\theta(J_{S_\theta(X)}) 
            = - \mathbb{E}_\theta \left( \left( \frac{\partial ^2 \log f_\theta(X)}{\partial \theta_i \partial \theta_j}\right)_{i,j}\right)
        $$
        avec $J_{S_\theta(X)}$ la Jacobienne du score. 
    \end{itemize}
\end{prop}

Ces propriétés seront donc privilégiées pour le calcul pratique de l'information de Fisher. 

\begin{prop}[Cas du modèle exponentiel canonique]
    Dans le cas d'un modèle exponentiel canonique, on a $ \mathcal{I}_\theta(X) = \V_\theta(T(X)) = - H_{ \kappa_\theta}$ 
    la hessienne de $ \kappa_\theta = \log k_\theta$. 
    Donc si $ p = 1$, alors $ \mathcal{I}_\theta(X) = - \frac{ \partial ^2}{ \partial \theta^2} \log k_\theta$. 
\end{prop}

\begin{proposition}
    Comme dit précédement, les modèles exponentiels de plein rang sont très utiles. 
    Ils permettent de calculer simplement l'information de Fisher à partir des fonctions 
    définissant la densité. On a donc : 
        \begin{align*}
            \mathcal{I}_\theta(X) &\underset{\text{def}}{=} J_{\gamma}^t \times \V_\theta(T(X)) \times J_\gamma(\theta) \\ 
            &\underset{\text{prop}}{=} - H_{ \kappa_\theta} - H_{\gamma} \left( J_{\gamma}^{-1}\right)^t \nabla_\theta \kappa_\theta
        \end{align*}
    Si $p = 1$, 
        $$ \mathcal{I}_\theta(X) = ( \gamma'(\theta))^2 \times V_\theta(T(X))
            = \frac{- k_\theta'}{k_\theta} \left( \frac{k_\theta ''}{k_\theta'} + \frac{ \gamma_\theta ''}{\gamma_\theta'} 
            - \frac{k_\theta'}{k_\theta}\right)
        $$
\end{proposition}

\begin{example}
    Soit le modèle $X \sim \mathcal{N}( \theta_1, \theta_2^2)$, avec 
    $ 
        \begin{cases}
            \theta_1 = \mathbb{E}_\theta(X) \\
            \theta_2 = \sqrt{\V_\theta(X)}
        \end{cases}
    $. 
    On a donc : 
        $$ f_\theta(x) = \frac{1}{2 \pi} \exp \left( \langle T(x), \gamma(\theta) \rangle + \kappa_\theta\right) $$
    où $ p =q = 2$ et : 
        $$ 
            T(x) = (x, x^2), \quad \kappa_\theta = - \frac{ \theta_1^2}{2 \theta_2^2} - \log \theta_2, \quad 
            \gamma(\theta) = \left( \frac{ \theta_1}{ \theta_2}, - \frac{1}{ 2 \theta_2^2}\right) 
        $$
    On peut donc ensuite appliquer les formules : 
        $$ J_\gamma = 
            \begin{pmatrix}
                \frac{1}{ \theta_2^2} & - \frac{2 \theta_1}{ \theta_2^2} \\
                0 & \frac{1}{ \theta_2^2}
            \end{pmatrix}
            \quad \V_\theta(T(X)) = 
            \begin{pmatrix}
                \theta_2^2 & Cov(X, X^2) \\ 
                Cov(X, X^2) & \V_\theta(X^2)
            \end{pmatrix}
        $$ 
    On a donc : 
        \begin{align*}
            \mathcal{I}_\theta(X) &= Y_\gamma^t \times V_\theta(T(X)) \times J_\gamma \\ 
            &= - \mathbb{E}_\theta \left( \frac{\partial^2}{ \partial \theta_i \partial \theta_j} 
                \left(
                    - \frac{(X- \theta_1)^2}{2 \theta_2^2} - \log \theta_2
                \right)
                \right)_{i,j}
            = 
                \begin{pmatrix}
                    \frac{1}{ \theta_2^2} & 0 \\ 
                    0 & \frac{2}{ \theta_2^2}
                \end{pmatrix}
        \end{align*}
\end{example}

% ===================================================================================================================
% Propriétés de l'information de Fisher

\section{Propriétés de l'information de Fischer}

\subsection{Additivité}

Les propriétés de l'information de Fisher découlent de celles du score défini 
à partir d'un $ \log$. Ainsi, l'information de Fisher va hériter de sa propriété 
d'additivité. 

\begin{proposition}
    Soient $S,T$ deux statistiques $ \mathcal{P}_\Theta$-indépendantes. Si $ \myP_\Theta^S$ et 
    $ \myP_\Theta^T$ vérifient $ (H_{1, 2, 3}) $, alors l'information du modèle $ \myP_\Theta^{S,T}$
    est $ \mathcal{I}_\theta(S, T) = \mathcal{I}_\theta(S) + \mathcal{I}_\theta(T)$. 
\end{proposition}

\begin{corollary}[Additivité]
    Dans le modèle $ \mathcal{P}_\Theta ^{\otimes n}$, l'information de Fisher sur $ \theta$ est : 
        $$ \mathcal{I}_\theta(X_1, \dots, X_n) = n \mathcal{I}_\theta(I_1). $$
\end{corollary}

\subsection{Décroissance}

L'information de Fisher donne la quentité d'information contenue dans l'acte d'observer 
la réalisation d'un échantillon dans un certain espace. Or les statistiques nous permettent 
de condenser cette inormation dans un autre espace. Déterminer le lien de cette \emph{information image}
par rapport à l'initiale. 

\begin{definition}[Information image]
    Soit $ S : ( \mathcal{X}, \mathcal{A}) \longrightarrow ( \mathcal{Y}, \mathcal{B})$ une statistique. 
    Si $( \mathcal{Y}, \mathcal{B}, \mathcal{P}_\Theta^S)$ vérifie $(H_{1,2})$ avec $ \mu^S$ une mesure 
    $ \sigma$-finie, alors $ \mathcal{I}_\theta(S)$ est l'information de Fisher dans ce modèle 
    sur $ \theta$ dans l'acte d'observer $ S$.     
\end{definition}

\begin{lemma}
    Si $ \mathcal{P}_\Theta$ et $ \mathcal{P}_\Theta^S$ vérifient $(H_{1, 2, 3})$, alors : 
        $$ S_\theta(S) = \nabla_\theta \log f_{s,\theta}(S) = \mathbb{E}_\theta( \nabla_\theta \log L_\theta(X) \mid S) 
        = \mathbb{E}_\theta(S_\theta(S) \mid S) 
        $$
    On a aussi : 
        $$ \mathcal{I}_\theta(S) = \V_\theta(S_\theta(S)) = \V_\theta( \mathbb{E}_\theta(S_\theta(X)) \mid S). $$
\end{lemma}

Ainsi, le score image $S_\theta(S)$ est l'espérance du score initial conditionnellement à $S$. 

\begin{prop}[Décroissance]
    Si $ \mathcal{P}_\Theta$ et $ \mathcal{P}_\Theta^S$ vérifient $(H_{1,2,3})$, alors 
    $ \mathcal{I}_\theta(X) - \mathcal{I}_\theta(S)$ est semi-définie positive. 
    De même, dans le cas du modèle d'échantillonnage,  $ n \mathcal{I}_\theta(X_1) - \mathcal{I}_\theta(S(X_1, \dots, X_n))$
    l'est aussi. 
\end{prop}

\subsection{Exhaustivité}

Tout l'intérêt des statistiques exhaustives réside dans la propriété suivante. 

\begin{proposition}
    Sous $(H_{1, 2, 3})$, si $S$ est exhaustive, il n'y a pas de perte d'information par transfert. 
        $$ i.e \quad n \mathcal{I}_\theta(X) = \mathcal{I}_\theta(S(X_1, \dots, X_n)) $$
\end{proposition}


\subsection{Contaste de Kullback}

\begin{proposition}
    Sous $(H_{1, 2, 3, 4})$, pour tout $ \theta_0 \in \Theta$ fixé, on a : 
        $$ \mathcal{I}_{\theta_0}(X) = \left( \left( \frac{ \partial^2 K( \theta_0 \mid \theta)}{ \partial \theta_i \partial \theta_j} \right)_{i,j}\right) $$
\end{proposition}



